% `tex.web` §1197: a display opened with `$$` has to close with `$$` too.
%
%     begin get_x_token;
%     if cur_cmd<>math_shift then
%       begin print_err("Display math should end with $$");
%       help2("The `$' that I just saw supposedly matches a previous `$$'.")
%       ("So I shall assume that you typed `$$' both times."); back_error;
%       end;
%     end
%
% We read the token with `get_next` and put back anything that was not a `$`
% in silence, so the commonest way to mistype a display — one `$` at the end
% instead of two — went by without a word, and the rest of the paragraph came
% out looking deliberate.
%
% Two details of where §1194 reaches this matter as much as the message:
%
%   * `get_x_token`, so a macro that *expands* to `$` closes the display and
%     raises nothing. `\let\dollar=$` does too — a `\let` to a character token
%     keeps its catcode.
%
%   * §1194 checks the fonts *first*, then `fin_mlist`, then this. So when a
%     formula is deleted for want of fonts the error context still shows the
%     second `$` unread, one token earlier than where we had it.
\catcode`\{=1 \catcode`\}=2 \catcode`\$=3 \catcode`\#=6
\hsize=100pt \parindent=0pt
\font\rm=cmr10 \textfont0=\rm \scriptfont0=\rm \scriptscriptfont0=\rm \rm
\font\mi=cmmi10 \textfont1=\mi \scriptfont1=\mi \scriptscriptfont1=\mi
\font\sy=cmsy10 \textfont2=\sy \scriptfont2=\sy \scriptscriptfont2=\sy
\font\ex=cmex10 \textfont3=\ex \scriptfont3=\ex \scriptscriptfont3=\ex

% One `$` where two belong. The token that was not a `$` goes back — the `C`
% is set, and the paragraph is the one the document asked for.
\setbox0\vbox{\noindent A$$B$C\par}
\message{PROBE:one=[\the\ht0]}

% ...and with a control sequence there instead, so the context carries the
% `<to be read again>` line above it.
\setbox1\vbox{\noindent A$$B$\relax C\par}
\message{PROBE:cs=[\the\ht1]}

% An expansion closes it: `get_x_token`, not `get_next`.
\def\closer{$}
\setbox2\vbox{\noindent A$$B$\closer C\par}
\message{PROBE:expanded=[\the\ht2]}

% ...and so does a `\let` alias, which is a character token and not a macro.
\let\dollar=$
\setbox3\vbox{\noindent A$$B$\dollar C\par}
\message{PROBE:let=[\the\ht3]}

% §1194's order: the font check comes first, so this reports the missing
% family with the display's second `$` still unread.
\textfont2=\nullfont \scriptfont2=\nullfont \scriptscriptfont2=\nullfont
\setbox4\vbox{\noindent A$$B$$C\par}
\message{PROBE:danger=[\the\ht4]}

% ...and both faults at once, in the order §1194 raises them.
\setbox5\vbox{\noindent A$$B$C\par}
\message{PROBE:both=[\the\ht5]}
\end
